The next proposition isolates the transfer operator in the one-step binary
case.  The following theorem is then obtained by averaging the deterministic
adjacent defect against this two-state specialization.

\begin{proposition}[Hidden tail factor and transfer spectrum]
\label{prop:hidden-tail-transfer}
Let $P=(p_{ab})_{a,b\in\{0,1\}}$ be an irreducible one-step transition matrix
on $\{0,1\}$.  For $a\in\{0,1\}$ and $\ell\ge 0$, let $Q_\ell^{(a)}$ be the
law of the tail statistic
\[
\beta(\omega_1\cdots\omega_\ell)
\]
when the Markov chain is started from the boundary state $\omega_0=a$.  Then
$Q_\ell^{(a)}$ is obtained from the four-state hidden chain
\[
Z_t\in\{E,O,Z_0,Z_1\},
\]
where $E$ and $O$ mean that the initial run of $1$'s is still alive with even
or odd length, while $Z_i$ is the absorbing output state reached after the
first $0$ when the final value of $\beta$ is $i$.  After the first symbol has
been read, the initial row vector is
\[
u_a=(0,p_{a1},p_{a0},0)
\]
in the ordered coordinates $(E,O,Z_0,Z_1)$, and all subsequent symbols are read
with transition matrix
\[
M=
\begin{pmatrix}
0 & p_{11} & p_{10} & 0\\
p_{11} & 0 & 0 & p_{10}\\
0 & 0 & 1 & 0\\
0 & 0 & 0 & 1
\end{pmatrix},
\]
with output vector
\[
h=(0,1,0,1)^T.
\]
Thus $Q_\ell^{(a)}(\{1\})=u_aM^{\ell-1}h$ for $\ell\ge 1$.  On the live
parity subspace the transfer operator is
\[
T=
\begin{pmatrix}
0 & p_{11}\\
p_{11} & 0
\end{pmatrix},
\]
whose eigenvalues are $p_{11}$ and $-p_{11}$.  Consequently
\[
Q_0^{(a)}=\delta_0
\]
and, for $\ell\ge 1$,
\[
Q_\ell^{(a)}
=
\left(1-\frac{p_{a1}(1-(-p_{11})^\ell)}{1+p_{11}}\right)\delta_0
+
\frac{p_{a1}(1-(-p_{11})^\ell)}{1+p_{11}}\delta_1.
\]
The observable $\beta$ therefore descends from the one-sided shift to this
four-state hidden Markov factor, with the only nontrivial oscillatory mode
equal to $-p_{11}$.
\end{proposition}

\begin{proof}
By Theorem~\ref{thm:tail-rigidity}, $\beta(q)$ is the parity of the initial
run of $1$'s in $q$.  If the first symbol is $0$, the output is already frozen
at $0$; if it is $1$, the live run has odd length.  Thereafter, while the next
symbol remains $1$, the parity toggles with probability $p_{11}$.  When the
first $0$ appears, the current parity is frozen as the final output.  This is
exactly the hidden-chain representation with initial vector $u_a$, transition
matrix $M$, and output vector $h$.  The live block of $M$ is the displayed
matrix $T$, so its eigenvalues are $p_{11}$ and $-p_{11}$.

It remains only to compute the mass assigned to the output state $1$.  For
$\ell=0$ the tail word is empty, so $Q_0^{(a)}=\delta_0$.  For $\ell\ge 1$,
summing the disjoint events that the initial run has odd length gives
\[
\PP[\beta(\omega_1\cdots\omega_\ell)=1\mid \omega_0=a]
=
\sum_{\substack{1\le r<\ell\\ r\ {\rm odd}}}
p_{a1}p_{11}^{r-1}p_{10}
+
\mathbf{1}_{\{\ell\ {\rm odd}\}}p_{a1}p_{11}^{\ell-1}.
\]
If $\ell=2s$, this equals
\[
\sum_{j=0}^{s-1}p_{a1}p_{11}^{2j}(1-p_{11})
=
\frac{p_{a1}(1-p_{11}^{2s})}{1+p_{11}}.
\]
If $\ell=2s+1$, it equals
\[
\sum_{j=0}^{s-1}p_{a1}p_{11}^{2j}(1-p_{11})
+p_{a1}p_{11}^{2s}
=
\frac{p_{a1}(1+p_{11}^{2s+1})}{1+p_{11}}.
\]
These two cases combine to
\[
\frac{p_{a1}(1-(-p_{11})^\ell)}{1+p_{11}},
\]
which is the coefficient of $\delta_1$ in $Q_\ell^{(a)}$.  The coefficient of
$\delta_0$ is its complement.
\end{proof}

\begin{theorem}[Stationary one-step Markov extension]
\label{thm:markov-defect-law}
Let $\PP_P$ be the stationary irreducible one-step Markov measure on
$\{0,1\}^{\NN}$ with transition matrix
\[
P=(p_{ab})_{a,b\in\{0,1\}}
\]
and stationary distribution $\pi=(\pi_0,\pi_1)$.  Let
$\omega^\infty=(\omega_1,\omega_2,\dots)$ denote the coordinate process, and
for $m\ge 1$ write
\[
r_{\infty\to m}(\omega^\infty):=(\omega_1,\dots,\omega_m)\in\Omega_m.
\]
Fix $m\ge 1$ and $\ell\ge 0$, and set $n:=m+\ell$.  For $a\in\{0,1\}$ define
probability measures on $\{0,1\}^m$ by
\[
\lambda_{m,a}(v)
:=
\PP_P\!\left[
\kappa_{m+1\to m}(r_{\infty\to m}(\omega^\infty),1)=v
\ \middle|\
\omega_m=a
\right],
\]
and define
\[
\rho_\ell^{(a)}
:=
\PP_P\!\left[
\beta(\omega_{m+1}\cdots\omega_{m+\ell})=1
\ \middle|\
\omega_m=a
\right].
\]
Then
\[
\Law_{\PP_P}\!\bigl(D_{n\to m}(r_{\infty\to n}(\omega^\infty))\bigr)
=
\sum_{a=0}^1
\pi_a\bigl((1-\rho_\ell^{(a)})\delta_0+\rho_\ell^{(a)}\lambda_{m,a}\bigr),
\]
and
\[
\rho_0^{(a)}=0,
\qquad
\rho_\ell^{(a)}
=
\frac{p_{a1}(1-(-p_{11})^\ell)}{1+p_{11}}
\qquad (\ell\ge 1).
\]
Consequently
\[
\PP_P\!\bigl[D_{n\to m}(r_{\infty\to n}(\omega^\infty))\ne 0\bigr]
=
\sum_{a=0}^1
\pi_a\,\rho_\ell^{(a)}\bigl(1-\lambda_{m,a}(0)\bigr),
\]
the defect law converges as $\ell\to\infty$ to
\[
\nu_{m,P}
:=
\sum_{a=0}^1
\pi_a\left(
\left(1-\frac{p_{a1}}{1+p_{11}}\right)\delta_0
\;+\;
\frac{p_{a1}}{1+p_{11}}\lambda_{m,a}
\right),
\]
and
\[
d_{\mathrm{TV}}
\!\left(
\Law_{\PP_P}\!\bigl(D_{m+\ell\to m}(r_{\infty\to m+\ell}(\omega^\infty))\bigr),
\nu_{m,P}
\right)
\le
\frac{\pi_1}{1+p_{11}}\,p_{11}^{\ell}.
\]
If $p_{01}=p_{11}=u$, then
\[
\rho_\ell^{(0)}=\rho_\ell^{(1)}
=
\frac{u(1-(-u)^\ell)}{1+u},
\]
and the uniform law is recovered at $u=\frac12$.
\end{theorem}

\begin{proof}
Set
\[
P_{m,\ell}
:=
\Law_{\PP_P}\!\bigl(D_{m+\ell\to m}(r_{\infty\to m+\ell}(\omega^\infty))\bigr).
\]
Because the chain is irreducible on two states, $\pi_a>0$ for $a=0,1$, so the
conditional laws $\lambda_{m,a}$ are well defined.
Fix $p\in\Omega_m$ and write $a:=p_m$.  By the Markov property, conditioned on
$r_{\infty\to m}(\omega^\infty)=p$, the future block
$q:=(\omega_{m+1},\dots,\omega_{m+\ell})$ is independent of the past and is
distributed as a length-$\ell$ Markov chain started from state $a$.
Theorem~\ref{thm:tail-collapse} therefore gives
\[
\Law_{\PP_P}\!\bigl(D_{m+\ell\to m}(p,q)\,\big|\,r_{\infty\to m}(\omega^\infty)=p\bigr)
=
(1-\rho_\ell^{(a)})\delta_0+\rho_\ell^{(a)}\delta_{\kappa_{m+1\to m}(p,1)}.
\]
Conditioning further on $\omega_m=a$ and averaging over prefixes with last
digit $a$ yields
\[
\Law_{\PP_P}\!\bigl(
D_{m+\ell\to m}(r_{\infty\to m+\ell}(\omega^\infty))
\ \big|\ \omega_m=a
\bigr)
=
(1-\rho_\ell^{(a)})\delta_0+\rho_\ell^{(a)}\lambda_{m,a}.
\]
Since $\PP_P[\omega_m=a]=\pi_a$ by stationarity, averaging over $a$ proves the
first displayed formula for $P_{m,\ell}$.  The formula for the nonzero defect
probability is immediate.

Proposition~\ref{prop:hidden-tail-transfer}, applied to the future chain
started from the boundary state $a=\omega_m$, gives
\[
\rho_0^{(a)}=0,
\qquad
\rho_\ell^{(a)}
=
\frac{p_{a1}(1-(-p_{11})^\ell)}{1+p_{11}}
\qquad(\ell\ge 1).
\]

Let
\[
\rho_\infty^{(a)}:=\frac{p_{a1}}{1+p_{11}}.
\]
Since the chain is irreducible, one has $p_{11}<1$.  Therefore
\[
\nu_{m,P}
=
\sum_{a=0}^1
\pi_a\bigl((1-\rho_\infty^{(a)})\delta_0+\rho_\infty^{(a)}\lambda_{m,a}\bigr),
\]
so
\[
P_{m,\ell}-\nu_{m,P}
=
\sum_{a=0}^1
\pi_a\bigl(\rho_\ell^{(a)}-\rho_\infty^{(a)}\bigr)(\lambda_{m,a}-\delta_0).
\]
Since $d_{\mathrm{TV}}(\lambda_{m,a},\delta_0)\le 1$,
\[
d_{\mathrm{TV}}(P_{m,\ell},\nu_{m,P})
\le
\sum_{a=0}^1 \pi_a\bigl|\rho_\ell^{(a)}-\rho_\infty^{(a)}\bigr|
=
\frac{p_{11}^{\ell}}{1+p_{11}}
\sum_{a=0}^1 \pi_a p_{a1}.
\]
Stationarity gives $\pi_0p_{01}+\pi_1p_{11}=\pi_1$, hence
\[
d_{\mathrm{TV}}(P_{m,\ell},\nu_{m,P})
\le
\frac{\pi_1}{1+p_{11}}\,p_{11}^{\ell}.
\]
If $p_{01}=p_{11}=u$, then the chain is i.i.d.\ Bernoulli$(u)$, so
$\rho_\ell^{(0)}=\rho_\ell^{(1)}$.  At $u=\frac12$, the stationary law has
$\pi_0=\pi_1=\frac12$, and the stationary mixture of the two conditional
prefix laws is the unconditional uniform prefix law.  Hence
\[
\sum_{a=0}^1 \pi_a\lambda_{m,a}=\lambda_m.
\]
Substituting this identity into the displayed law for $P_{m,\ell}$ gives
\[
P_{m,\ell}
=(1-\rho_\ell)\delta_0+\rho_\ell\lambda_m,
\qquad
\rho_\ell=\frac{2^\ell-(-1)^\ell}{3\cdot 2^\ell},
\]
which is Theorem~\ref{thm:tail-collapse}.
\end{proof}

